\src{178}
\begin{align*}
\Extsh_{\Osh_X}^{p}(F,L)
&=L\otimes(\Omega_X^{n})'\otimes
  \Extsh_{\Osh_X}^{p}(F,\Omega_X^{n})\\
&=L\otimes(\Omega_X^{n})'\otimes
  \Homsh_{\Osh_X}(F,\omega_Y^{q})\\
&=\Homsh_{\Osh_X}\!\left(
  F\otimes L'\otimes\Omega_X^{n},\omega_Y^{q}\right),
\end{align*}
where $q=n-p$, which gives formula (4.1 bis).

In particular, take $F=\Osh_Y$ and $L=\Omega_X^{p}$. Since
$\Omega_X^{n}\otimes(\Omega_X^{p})'=\Omega_X^{n-p}$, one obtains a canonical
isomorphism
\begin{equation}\tag{4.2}
  \operatorname{Ext}_{\Osh_X}^{p}(X;\Osh_Y,\Omega_X^{p})
  =\operatorname{Hom}_{\Osh_X}
  (X;\Omega_X^{n-p},\omega_Y^{n-p}).
\end{equation}
Now suppose for simplicity that $Y$ is nonsingular, so that
$\omega_Y^{n-p}=\Omega_Y^{n-p}$. There is a natural homomorphism from
$\Omega_X^{n-p}$ to $\Omega_Y^{n-p}$, hence a canonical section $s_Y$ of the
sheaf $\Extsh_{\Osh_X}^{p}(\Osh_Y,\Omega_X^{p})$. If all components of $Y$ have
dimension $n-p$, we call it the \emph{fundamental section} of the sheaf
$\Extsh_{\Osh_X}^{p}(\Osh_Y,\Omega_X^{p})$. By (4.1), this section defines an
element of $\operatorname{Ext}_{\Osh_X}^{p}(X;\Osh_Y,\Omega_X^{p})$. The natural
homomorphism $\Osh_X\longrightarrow\Osh_Y$ in turn defines a homomorphism
\[
  \operatorname{Ext}_{\Osh_X}^{p}(X;\Osh_Y,\Omega_X^{p})
  \longrightarrow
  \operatorname{Ext}_{\Osh_X}^{p}(X;\Osh_X,\Omega_X^{p})
  =H^{p}(X,\Omega_X^{p}).
\]
Thus one obtains an element of $H^{p}(X,\Omega_X^{p})$, denoted $P_X(Y)$ and
called the \emph{cohomology class of $Y$ in $X$}. It is therefore obtained from
the section $s_Y$ of $\Extsh_{\Osh_X}^{p}(\Osh_Y,\Omega_X^{p})$ through the
following diagram of homomorphisms:
\begin{equation}\tag{4.3}
\begin{aligned}
H^{p}(X,\Omega_X^{p})
&=\operatorname{Ext}_{\Osh_X}^{p}(X;\Osh_X,\Omega_X^{p})\\
&\longleftarrow\operatorname{Ext}_{\Osh_X}^{p}(X;\Osh_Y,\Omega_X^{p})
\xrightarrow{\ \sim\ }
H^{0}\!\left(X,\Extsh_{\Osh_X}^{p}(\Osh_Y,\Omega_X^{p})\right).
\end{aligned}
\end{equation}

Call a \emph{nonsingular cycle of dimension $n-p$} an element of the free
abelian group generated by the irreducible nonsingular subvarieties of dimension
$n-p$ in $X$. Then the function $Y\longmapsto P(Y)$ extends to a homomorphism
from the group of nonsingular cycles of dimension $n-p$ on $X$ to the group
$H^{p}(X,\Omega_X^{p})$.

Let $Z^{n-p}$ and $Z'^{n-p'}$ be nonsingular cycles of dimensions $n-p$ and
$n-p'$. They are said to \emph{intersect transversely} if every component of
$Z$ intersects every component of $Z'$ transversely. The cycle $Z.Z'$ is then
defined; it is a nonsingular cycle of dimension $n-p-p'$. With this understood,
we have the

